\documentclass[aspectratio=169,10pt]{beamer}
\usetheme{Madrid}
\usecolortheme{default}
\usefonttheme{professionalfonts}

\usepackage[T1]{fontenc}
\usepackage[utf8]{inputenc}
\usepackage[french]{babel}
\usepackage{lmodern}
\usepackage{tikz}
\usetikzlibrary{positioning,arrows.meta,shapes.geometric,fit,calc}
\usepackage{booktabs}
\usepackage{listings}
\usepackage{xcolor}

\definecolor{Primary}{HTML}{0D3B66}
\definecolor{Accent}{HTML}{EE964B}
\definecolor{Soft}{HTML}{F4D35E}
\setbeamercolor{structure}{fg=Primary}
\setbeamercolor{title}{fg=Primary}
\setbeamercolor{frametitle}{fg=white,bg=Primary}
\setbeamercolor{block title}{bg=Primary!12, fg=Primary}
\setbeamercolor{block body}{bg=Primary!4, fg=black}
\setbeamercolor{itemize item}{fg=Primary}
\setbeamercolor{itemize subitem}{fg=Primary}
\setbeamertemplate{headline}{}

\setbeamertemplate{frametitle}{
  \nointerlineskip
  \begin{beamercolorbox}[wd=\paperwidth,ht=3.2ex,dp=1.2ex,leftskip=1em,rightskip=1em]{frametitle}
    \usebeamerfont{frametitle}\insertframetitle
  \end{beamercolorbox}
}

\setbeamertemplate{navigation symbols}{}

\lstset{
  basicstyle=\ttfamily\scriptsize,
  breaklines=true,
  columns=fullflexible,
  frame=single,
  rulecolor=\color{Primary!40},
  keywordstyle=\color{Primary}\bfseries,
  commentstyle=\color{gray},
  showstringspaces=false
}

\title[Plateforme IA multi-agents]{Plateforme IA de livraison multi-agents\\\large Du besoin aux artefacts logiciels en quelques minutes}
\subtitle{Pitch Hackathon -- 12 min + 8 min Q\&A}
\author{Équipe 24\\Ethan Brosselard \and Savain Blaneus \and Rayane Malik}
\institute{Challenge IBM}
\date{Mars 2026}

\begin{document}

% 1) Titre
\section{Introduction}
\begin{frame}{Plateforme IA de livraison multi-agents}
  \vspace{0.4cm}
  \begin{center}
    {\Large\textbf{Du besoin client au prototype démontrable}}\\[0.3cm]
    {\normalsize Pitch hackathon orienté jury entreprise}\\[0.8cm]
  \end{center}

  \begin{block}{Ce que vous allez voir}
  \begin{itemize}
    \item Un problème réel de delivery logiciel
    \item Une solution multi-agents concrète
    \item Une démo reproductible et traçable
  \end{itemize}
  \end{block}

  \vfill
  \begin{center}
    {\small Équipe 24 -- Ethan Brosselard, Savain Blaneus, Rayane Malik}\\
    {\small Challenge IBM -- Mars 2026}
  \end{center}
\end{frame}

% Sommaire
\section{Sommaire}
\begin{frame}{Sommaire\\Plan du pitch}
\tableofcontents
\end{frame}

% 2) Accroche
\section{Contexte et problème}
\begin{frame}{Imaginez cette scène...}
\begin{columns}[T]
\column{0.62\textwidth}
\begin{block}{17h12, veille de démo client}
\begin{itemize}
  \item Le client dit : \textbf{« Je veux un prototype demain matin. »}
  \item L'équipe a : un brief, des notes, 3 idées contradictoires.
  \item En quelques heures, il faut aligner \textbf{métier + architecture + code}.
\end{itemize}
\end{block}
\vspace{0.15cm}
\begin{alertblock}{La vraie douleur}
Ce n'est pas seulement coder vite.\\
C'est livrer vite \textbf{sans perdre la cohérence} ni la confiance du client.
\end{alertblock}
\column{0.38\textwidth}
\begin{tikzpicture}[node distance=5mm]
\node[draw,rounded corners,fill=red!10,align=center,text width=3.4cm,minimum height=0.85cm] (a) {17h : brief incomplet};
\node[draw,rounded corners,fill=red!10,below=of a,align=center,text width=3.4cm,minimum height=0.85cm] (b) {20h : décisions floues};
\node[draw,rounded corners,fill=red!10,below=of b,align=center,text width=3.4cm,minimum height=0.85cm] (c) {23h : dette de qualité};
\node[draw,rounded corners,fill=green!14,below=of c,align=center,text width=3.4cm,minimum height=1.0cm] (d) {8h : démo crédible};
\draw[-{Latex[length=2mm]},thick] (a)--(b);
\draw[-{Latex[length=2mm]},thick] (b)--(c);
\draw[-{Latex[length=2mm]},thick,Primary] (c)--(d);
\end{tikzpicture}
\end{columns}
\end{frame}

% 3) Problème
\begin{frame}{Le problème concret}
\begin{columns}[T]
\column{0.52\textwidth}
\begin{block}{Situation de départ}
\textbf{Besoin métier} non structuré\\
$\Downarrow$\\
\textbf{Livrables attendus} en 24h
\end{block}
\vspace{0.15cm}
\begin{alertblock}{Friction}
Spécification \(\neq\) Architecture \(\neq\) Code
\end{alertblock}
\column{0.48\textwidth}
\begin{tikzpicture}[node distance=4.5mm, every node/.style={font=\scriptsize}]
\node[draw,rounded corners,fill=red!10,text width=3.6cm,align=center,minimum height=0.8cm] (m) {MÉTIER\\brief + urgences};
\node[draw,rounded corners,fill=red!10,below=of m,text width=3.6cm,align=center,minimum height=0.8cm] (a) {ARCHI\\décisions tardives};
\node[draw,rounded corners,fill=red!10,below=of a,text width=3.6cm,align=center,minimum height=0.8cm] (c) {CODE\\hétérogène};
\node[draw,rounded corners,fill=Accent!20,below=of c,text width=3.6cm,align=center,minimum height=0.9cm] (r) {RÉSULTAT\\démo fragile};
\draw[-{Latex[length=2mm]},thick] (m)--(a);
\draw[-{Latex[length=2mm]},thick] (a)--(c);
\draw[-{Latex[length=2mm]},thick] (c)--(r);
\end{tikzpicture}
\end{columns}
\end{frame}

% 4) Enjeux
\section{Enjeux}
\begin{frame}{Pourquoi c'est critique ?}
\begin{columns}[T]
\column{0.5\textwidth}
\textbf{Enjeux business}
\begin{itemize}
  \item Délai de prototypage
  \item Coût de coordination
  \item Qualité perçue en démo
\end{itemize}
\column{0.5\textwidth}
\textbf{Enjeux techniques}
\begin{itemize}
  \item Cohérence inter-artefacts
  \item Gouvernance des décisions
  \item Reproductibilité des runs
\end{itemize}
\end{columns}
\vspace{0.4cm}
\begin{alertblock}{Contrainte hackathon IBM}
Démontrer une solution multi-agents \textbf{défendable en entreprise} : traçable, lisible, fiable.
\end{alertblock}
\end{frame}

% 5) Solution proposée
\section{Solution}
\begin{frame}{Notre solution}
\begin{block}{Plateforme IA de livraison multi-agents}
Un \textbf{générateur Python multi-agents} qui transforme une exigence en artefacts clés d'une chaîne logicielle.
\end{block}
\vspace{0.2cm}
\begin{itemize}
  \item Entrées : texte, Markdown, JSON.
  \item Sorties : analyse, plan, C4, code app, tests, CI/CD, traces Plan/Act/Reason.
  \item Exécution via CLI \textbf{et} API.
\end{itemize}
\vspace{0.2cm}
\begin{center}
\textbf{Objectif :} industrialiser la phase \og de l'idée au prototype crédible \fg.
\end{center}
\end{frame}

% 6) Idée principale
\section{Système et architecture}
\begin{frame}{Idée principale du système}
\begin{center}
\begin{tikzpicture}[node distance=6mm, every node/.style={font=\small}]
\node[draw,rounded corners,fill=Primary!8,minimum width=2.8cm,minimum height=0.8cm] (in) {Input besoin};
\node[draw,rounded corners,fill=Primary!8,right=8mm of in,minimum width=2.8cm,minimum height=0.8cm] (state) {État typé central};
\node[draw,rounded corners,fill=Primary!8,right=8mm of state,minimum width=2.8cm,minimum height=0.8cm] (agents) {Agents spécialisés};
\node[draw,rounded corners,fill=Accent!16,right=8mm of agents,minimum width=2.8cm,minimum height=0.8cm] (out) {Artefacts prêts};
\draw[-{Latex[length=2mm]},thick] (in)--(state);
\draw[-{Latex[length=2mm]},thick] (state)--(agents);
\draw[-{Latex[length=2mm]},thick] (agents)--(out);
\end{tikzpicture}
\end{center}
\vspace{0.3cm}
\begin{itemize}
  \item \textbf{Décomposition des rôles} pour réduire l'ambiguïté.
  \item \textbf{Flux déterministe} pour fiabilité démo.
  \item \textbf{Contrats Pydantic} pour valider chaque étape.
\end{itemize}
\end{frame}

% 7) Architecture
\begin{frame}{Architecture système (vue conteneurs)}
\begin{center}
\begin{tikzpicture}[node distance=7mm, every node/.style={font=\scriptsize}]
\node[draw,rounded corners,fill=blue!8,minimum width=3.1cm,minimum height=0.9cm] (cli) {CLI / API FastAPI};
\node[draw,rounded corners,fill=blue!8,right=of cli,minimum width=3.2cm,minimum height=0.9cm] (graph) {Orchestrateur LangGraph};
\node[draw,rounded corners,fill=blue!8,right=of graph,minimum width=3.2cm,minimum height=0.9cm] (llm) {Provider LLM abstrait};

\node[draw,rounded corners,fill=green!8,below=of graph,minimum width=10.4cm,minimum height=0.9cm] (agents) {Agents: spec\_analyst $\rightarrow$ architect $\rightarrow$ developer $\rightarrow$ qa\_devops $\rightarrow$ reviewer};
\node[draw,rounded corners,fill=orange!10,below=of agents,minimum width=10.4cm,minimum height=1.0cm] (art) {Artifacts: workspace/generated\_app + outputs/plans, outputs/c4, outputs/traces};

\draw[-{Latex[length=2mm]},thick] (cli)--(graph);
\draw[-{Latex[length=2mm]},thick] (graph)--(llm);
\draw[-{Latex[length=2mm]},thick] (graph)--(agents);
\draw[-{Latex[length=2mm]},thick] (agents)--(art);
\end{tikzpicture}
\end{center}
\end{frame}

% 8) Fonctionnement global
\section{Fonctionnement}
\begin{frame}{Fonctionnement global (pipeline)}
\begin{center}
\begin{tikzpicture}[node distance=6mm, every node/.style={font=\scriptsize}]
\node[draw,rounded corners,fill=Primary!10,text width=2.2cm,align=center,minimum height=0.9cm] (s1) {1. Analyse\\besoin};
\node[draw,rounded corners,fill=Primary!10,right=5mm of s1,text width=2.2cm,align=center,minimum height=0.9cm] (s2) {2. Plan\\+ Archi};
\node[draw,rounded corners,fill=Primary!10,right=5mm of s2,text width=2.2cm,align=center,minimum height=0.9cm] (s3) {3. Génération\\App};
\node[draw,rounded corners,fill=Primary!10,right=5mm of s3,text width=2.2cm,align=center,minimum height=0.9cm] (s4) {4. C4 + CI\\+ Tests};
\node[draw,rounded corners,fill=Accent!20,right=5mm of s4,text width=2.2cm,align=center,minimum height=0.9cm] (s5) {5. Rapport\\+ Traces};
\draw[-{Latex[length=2mm]},thick] (s1)--(s2);
\draw[-{Latex[length=2mm]},thick] (s2)--(s3);
\draw[-{Latex[length=2mm]},thick] (s3)--(s4);
\draw[-{Latex[length=2mm]},thick] (s4)--(s5);
\end{tikzpicture}
\end{center}
\vspace{0.2cm}
\begin{alertblock}{Supervision humaine}
\textbf{Pause $\rightarrow$ intervention $\rightarrow$ reprise} aux points critiques.
\end{alertblock}
\end{frame}

% 9) Démonstration concrète
\section{Démonstration}
\begin{frame}{Démonstration (ce que voit le jury)}
\begin{columns}[T]
\column{0.50\textwidth}
\textbf{Parcours live (90 s)}
\vspace{0.15cm}
\begin{enumerate}
  \item \textbf{On lance} le run\\
  \scriptsize{\texttt{ai\_delivery.cli run --input inputs/sample\_spec.md}}
  \item \textbf{On suit} les étapes agents\\
  \scriptsize{spec\_analyst $\rightarrow$ architect $\rightarrow$ developer...}
  \item \textbf{On ouvre} les artefacts générés\\
  \scriptsize{\texttt{workspace/generated\_app} + \texttt{outputs/*}}
\end{enumerate}
\vspace{0.15cm}
\textbf{Pilotage API}
\begin{itemize}
  \item \scriptsize \texttt{POST /runs}
  \item \scriptsize \texttt{POST /runs/\{id\}/resume}
  \item \scriptsize \texttt{POST /runs/\{id\}/interventions}
\end{itemize}
\column{0.50\textwidth}
\begin{tikzpicture}[node distance=5.5mm, every node/.style={font=\scriptsize}]
\node[draw,rounded corners,fill=Primary!10,text width=3.8cm,minimum height=0.9cm,align=center] (n1) {1) ENTRÉE\\spécification};
\node[draw,rounded corners,fill=Primary!10,right=9mm of n1,text width=3.8cm,minimum height=0.9cm,align=center] (n2) {2) EXÉCUTION\\agents spécialisés};
\node[draw,rounded corners,fill=Primary!10,below=9mm of n2,text width=3.8cm,minimum height=0.9cm,align=center] (n3) {3) CONTRÔLE\\pause / intervention};
\node[draw,rounded corners,fill=Accent!20,left=9mm of n3,text width=3.8cm,minimum height=0.9cm,align=center] (n4) {4) SORTIES\\app + C4 + traces};
\draw[-{Latex[length=2mm]},thick] (n1)--(n2);
\draw[-{Latex[length=2mm]},thick] (n2)--(n3);
\draw[-{Latex[length=2mm]},thick] (n3)--(n4);
\draw[-{Latex[length=2mm]},thick,dashed] (n4) to[bend left=25] node[midway,below,font=\tiny]{itération rapide} (n1);
\end{tikzpicture}
\end{columns}
\end{frame}

% 10) Résultats
\section{Résultats}
\begin{frame}{Résultats obtenus}
\begin{center}
\begin{tikzpicture}[node distance=6mm, every node/.style={font=\scriptsize}]
\node[draw,rounded corners,fill=Primary!10,text width=2.6cm,minimum height=1.1cm,align=center] (k1) {CLI + API\\opérationnels};
\node[draw,rounded corners,fill=Primary!10,right=6mm of k1,text width=2.6cm,minimum height=1.1cm,align=center] (k2) {App générée\\workspace/generated\_app};
\node[draw,rounded corners,fill=Primary!10,right=6mm of k2,text width=2.6cm,minimum height=1.1cm,align=center] (k3) {Docs C4\\context/container/component};
\node[draw,rounded corners,fill=Primary!10,right=6mm of k3,text width=2.6cm,minimum height=1.1cm,align=center] (k4) {Traces + plans\\auditables};
\end{tikzpicture}
\end{center}
\vspace{0.25cm}
\begin{block}{Message clé}
\textbf{Un run = une base projet cohérente, démontrable et industrialisable.}
\end{block}
\end{frame}

% 11) Valeur ajoutée
\section{Valeur}
\begin{frame}{Valeur ajoutée du projet}
\begin{columns}[T]
\column{0.5\textwidth}
\textbf{Pour le métier}
\begin{itemize}
  \item Prototype plus rapide
  \item Plus de lisibilité
  \item Meilleure capacité de décision
\end{itemize}
\column{0.5\textwidth}
\textbf{Pour la tech}
\begin{itemize}
  \item Stack standard (Python/FastAPI/React)
  \item Architecture modulaire
  \item Traçabilité native
\end{itemize}
\end{columns}
\vspace{0.3cm}
\begin{center}
\fcolorbox{Primary}{Primary!7}{\parbox{0.82\linewidth}{\centering
\textbf{Ce n'est pas seulement du code généré :} c'est une \textbf{chaîne de delivery complète} et auditable.
}}
\end{center}
\end{frame}

% 12) Crédibilité/faisabilité
\section{Crédibilité}
\begin{frame}{Crédibilité et faisabilité}
\begin{columns}[T]
\column{0.5\textwidth}
\textbf{Pourquoi c'est crédible ?}
\begin{itemize}
  \item Stack éprouvée
  \item Flux déterministe
  \item Contrats typés
  \item Traçabilité native
\end{itemize}
\column{0.5\textwidth}
\textbf{Pourquoi c'est faisable ?}
\begin{itemize}
  \item Rôles d'agents clairs
  \item Templates versionnés
  \item CLI/API déjà utilisables
  \item Extensible (Snowflake)
\end{itemize}
\end{columns}
\vspace{0.2cm}
\begin{alertblock}{Preuve de sérieux}
Le système produit des artefacts vérifiables, pas seulement un discours.
\end{alertblock}
\end{frame}

% 13) Limites
\section{Limites}
\begin{frame}{Limites actuelles (assumées)}
\begin{itemize}
  \item Intégration Snowflake réelle encore à finaliser.
  \item Backend généré encore partiellement en mémoire sur certains templates.
  \item Connexion frontend-backend à durcir pour une démo production-like.
  \item CI de l'application générée à renforcer.
\end{itemize}
\vspace{0.2cm}
\begin{alertblock}{Transparence}
Nous préférons expliciter les limites plutôt que sur-promettre.
\end{alertblock}
\end{frame}

% 14) Améliorations futures
\section{Roadmap}
\begin{frame}{Roadmap d'amélioration}
\begin{columns}[T]
\column{0.33\textwidth}
\textbf{P0}
\begin{itemize}
  \item Provider Snowflake réel
  \item SQLite end-to-end
\end{itemize}
\column{0.33\textwidth}
\textbf{P1}
\begin{itemize}
  \item UX intervention guidée
  \item CI générée robuste
\end{itemize}
\column{0.33\textwidth}
\textbf{P2}
\begin{itemize}
  \item Observabilité avancée
  \item API durcie (auth)
\end{itemize}
\end{columns}
\vspace{0.3cm}
\begin{block}{Vision}
Passer d'un \textbf{générateur de démo fiable} à un \textbf{accélérateur de delivery industriel}.
\end{block}
\end{frame}

% 15) Valeur entreprise
\section{Valeur entreprise}
\begin{frame}{Pourquoi ce projet a de la valeur pour une entreprise ?}
\begin{itemize}
  \item Réduction du time-to-prototype
  \item Standardisation des architectures
  \item Traçabilité des décisions
  \item Base projet plus propre
  \item Meilleure démo client
  \item Accélération innovation interne
\end{itemize}
\end{frame}

% 16) Conclusion mémorable
\section{Conclusion}
\begin{frame}{Conclusion -- message clé}
\begin{center}
{\Large \textbf{Demain matin, le client ne voit pas une promesse.\\Il voit une solution qui tourne.}}
\end{center}
\vspace{0.4cm}
\begin{itemize}
  \item Nous transformons un besoin flou en livrables concrets et alignés.
  \item Nous accélérons la delivery sans sacrifier gouvernance ni traçabilité.
  \item Nous donnons à l'entreprise une base propre pour innover plus vite.
\end{itemize}
\vspace{0.3cm}
\begin{center}
\fcolorbox{Primary}{Primary!8}{\parbox{0.83\linewidth}{\centering
\textbf{Plateforme IA de livraison multi-agents} = \\textit{accélérer sans perdre le contrôle, et convaincre plus vite.}
}}
\end{center}
\end{frame}

% 17) Q&A
\section{Questions}
\begin{frame}[standout]
\centering
{\Huge Questions ?}\\[0.4cm]
{\Large Merci !}
\vfill
\small 8 minutes de Q\&A
\end{frame}

\end{document}
